\documentclass[10pt,journal,compsoc]{IEEEtran}
% 補足資料（TAC の規則：付録は本文に入れず，別ファイルの補足資料として出す）．
\usepackage[T1]{fontenc}
\usepackage[utf8]{inputenc}
\usepackage{amsmath,amssymb}
\usepackage{graphicx}
\usepackage{booktabs}
\usepackage[nocompress]{cite}
\usepackage{url}
\usepackage[hidelinks]{hyperref}
% 自動生成．手で編集しない．
% 生成元: results/numbers.json ／ 生成器: scripts/gen_numbers_tex.py
% 本文の実験由来の数値は \NUM{key} で書く．未登録の key はビルド時にエラーになる．
\makeatletter
\providecommand{\NUM}[1]{%
  \@ifundefined{NUM@#1}{%
    \GenericError{}{numbers.tex: key '#1' is not registered in results/numbers.json}{}{}%
    \textbf{??}%
  }{\csname NUM@#1\endcsname}%
}
\makeatother

\newcommand{\FIG}[2]{\IfFileExists{figures/#1.pdf}{\includegraphics[width=#2]{figures/#1.pdf}}{\fbox{\parbox[c][6em][c]{\dimexpr#2-2\fboxsep-2\fboxrule\relax}{\centering\footnotesize figure \texttt{#1} (pending)}}}}
\renewcommand{\thesection}{S\arabic{section}}
\renewcommand{\thetable}{S\arabic{table}}
\renewcommand{\thefigure}{S\arabic{figure}}
\renewcommand{\theequation}{S\arabic{equation}}

\begin{document}
\title{Supplementary Material for ``What Few Target Labels Fix First: Threshold-versus-Direction Analysis and Cross-Modal Weak-Label Transfer for Cross-Corpus Multimodal Emotion Recognition''}
\author{Author~Names}
\maketitle

\section{Feature Extraction}
\label{sec:supp-features}

\emph{Face.}
From each video we take every $s$-th frame with $s = \lfloor 5f/25 \rfloor$ for a frame rate of $f$ frames per second, which gives five to six frames per second, and at most $400$ frames per utterance, sampled evenly.
IEMOCAP videos show two speakers side by side; we keep the half that shows the speaker.
YOLOv8n-face~\cite{yolov8face} detects faces with a confidence threshold of $0.25$, and we keep the largest face of each frame and resize it to $224\times224$ pixels.
The ResNet-50 of Ryumina et al.~\cite{ryumina2022facial}, trained on VGGFace2~\cite{cao2018vggface2} and fine-tuned on AffectNet~\cite{mollahosseini2019affectnet}, maps each crop to $512$ dimensions; the second encoder, the image encoder of CLIP ViT-L/14~\cite{radford2021clip}, maps the same crops (stored as boxes) to $1024$ dimensions.
Frame embeddings are averaged within segments of $20$ face frames shifted by $10$ frames, a final segment shorter than $10$ frames is dropped, and the segment means are averaged.
We do not use the standard deviation across segments, because utterances with fewer than $20$ face frames have one segment and a zero deviation, so that feature encodes utterance length.

\emph{Audio.}
Waveforms are resampled to $16$~kHz mono, truncated at $60$~s, normalized to zero mean and unit variance, and passed through WavLM-Large~\cite{chen2022wavlm} or HuBERT-Large~\cite{hsu2021hubert}; the final layer is averaged over time without padding ($1024$ dimensions).

\emph{Text.}
Transcripts are repaired for character-encoding errors, truncated at $256$ tokens, and passed through RoBERTa-base~\cite{liu2019roberta} or DeBERTa-v3-base~\cite{he2023debertav3}; the final layer is averaged over tokens ($768$ dimensions).

\emph{Counts.}
After filtering, the audio features contain \NUM{data-n-total-wavlm} utterances, the face features \NUM{data-n-total-resnet}, and the text features \NUM{data-n-total-roberta}.
Fusing audio and face keeps \NUM{data-n-impro-av}, \NUM{data-n-script-av}, \NUM{data-n-mosei-av}, \NUM{data-n-meld-av}, \NUM{data-n-afew-av}, and \NUM{data-n-cremad-av} utterances for the improvised and scripted IEMOCAP sessions, CMU-MOSEI, MELD, AFEW, and CREMA-D, and fusing all three modalities keeps \NUM{data-n-impro-avt}, \NUM{data-n-script-avt}, \NUM{data-n-mosei-avt}, and \NUM{data-n-meld-avt} utterances in the four groups with text.

\section{Derivations}
\label{sec:supp-derivations}

\emph{Threshold cost.}
Fix the direction of a pairwise decision and let the score difference $z$ be Gaussian with means $\mu_k$, $\mu_l$ and common standard deviation $\sigma$.
Balancing the two emotions places the threshold at $(\mu_k + \mu_l)/2$.
Estimated from $K$ labels per emotion, the two class means have variance $\sigma^2/K$ each, so the midpoint has variance $\sigma^2/(2K)$ and the residual offset satisfies $\mathbb{E}[o^2] = 1/(2K)$.
Since $\partial a/\partial o = 0$ at $o=0$ and $\partial^2 a/\partial o^2 = -(d/2)\,\phi(d/2)$ there,
\begin{equation}
a(d, 0) - a(d, o) \approx \frac{d}{4}\,\phi\!\left(\frac{d}{2}\right) o^2,
\end{equation}
and taking the expectation gives the threshold cost of the main text. The factor $(d/4)\phi(d/2)$ is largest at $d=2$, where it equals $\phi(1)/2$.

\emph{Direction cost.}
For a linear rule estimated from $K$ labels per emotion, the angle $\theta$ between the estimated and the best direction satisfies $\tan^2\theta \approx q/K$ to first order, where $q$ is an effective dimension; for the sample linear discriminant with $p$ features and squared Mahalanobis distance $\Delta^2$, the first-order term of the expansion of the misclassification rate gives $q = 2(p-1)/\Delta^2$~\cite{okamoto1963asymptotic,raudys1998intrinsic}.
The separation along the estimated direction is $\Delta\cos\theta = \Delta/\sqrt{1 + \tan^2\theta}$, which gives $d(K) = \Delta/\sqrt{1 + q/K}$.
With an $\ell_2$ penalty the effective dimension is smaller than $p$, and we fit $q$ to target-only learning curves.

\emph{Value of a source.}
Write the noise-to-signal ratio of a direction estimate as $\tan^2\theta$.
The target estimate from $K$ labels has ratio $q/K$, and the source direction, with efficiency $e = \cos\theta_s$, has ratio $(1-e^2)/e^2$.
Two independent estimates combined with inverse-variance weights have the ratio $\left[K/q + e^2/(1-e^2)\right]^{-1}$, which equals the ratio of a target-only estimate from $K + n$ labels with $n = q e^2/(1-e^2)$.

\emph{Correction of the threshold by pooling.}
Write the pooled objective as the source-only objective plus the loss of $N_t$ target points, and take one Newton step from the source-only solution.
With the curvature of the loss approximated by its mean $\kappa$ per domain, the step changes the score at the target mean by $\rho/(1+\rho)$ times the source-only offset, with $\rho$ given by the expression of the main text; the term $1/(N_s\kappa_s)$ comes from the intercept and the term $\boldsymbol{\delta}^\top(A_s + I/C)^{-1}\boldsymbol{\delta}$ from the weights.
The approximation keeps the loss of the target points to second order in the offset, takes the direction from the source, and omits the scatter of the target points themselves.
Fig.~\ref{fig:supp-lambda} compares the expression with the measured correction on real corpus pairs; the expression orders the pairs well but overestimates the average correction.

\section{Additional Results for Part~I}
\label{sec:supp-part1}

\begin{figure}[t]
\centering
\FIG{lambda}{\columnwidth}
\caption{Predicted against measured threshold correction $\lambda$ of pooled classifiers, per ordered pair and $K \in \{1, 5, 10\}$.}
\label{fig:supp-lambda}
\end{figure}

\begin{figure}[t]
\centering
\FIG{centering}{\columnwidth}
\caption{Source-only classifiers with and without centering, refitting only the intercepts on $K=5$ labels per emotion, and pooling with the same labels, for WavLM audio, HuBERT audio, and the face (three iterations). Without centering the pooled classifier takes the $C$ of the source-only classifier; with centering, $C$ is selected by cross-validation.}
\label{fig:supp-centering}
\end{figure}

\begin{figure}[t]
\centering
\FIG{sourcevalue}{\columnwidth}
\caption{What a source is worth. Predicted $n = q e^2/(1-e^2)$ against $n$ measured on the target-only learning curve, in labels per emotion, for single sources and all sources pooled ($K=5$ and $10$; without the two pairs of IEMOCAP sessions). Points beyond the learning curve are in the top band, and points with $n \le 0$ in the bottom band.}
\label{fig:supp-sourcevalue}
\end{figure}

\section{Gain Map with Cross-Validated Penalty}
\label{sec:supp-cvmap}

\begin{figure*}[t]
\centering
\FIG{cvmap}{\textwidth}
\caption{Gain of pooling when the penalty $C$ of the pooled classifier is selected by cross-validation on the $K$ target labels (solid) and when it is fixed to the $C$ of the source-only classifier as in the main text (dotted), averaged over ordered pairs. (a)~A, A~(HuBERT), V, and AV with thirty iterations. (b)~Configurations with text with five iterations.}
\label{fig:supp-cvmap}
\end{figure*}

Selecting $C$ on the target labels lowers the gain at small $K$, partly through ties at $K=1$ that the tie rule resolves toward a weak penalty (Section~3 of the main text).

\section{Additional Results for Part~II}
\label{sec:supp-part2}

\emph{Values of the start.}
The five values of Section~5.1 of the main text take three candidates each: $\gamma \in \{-0.5, 0, 0.5\}$, $\rho_s \in \{0.3, 0.5, 0.7\}$, $C_b \in \{2, 8, \infty\}$ ($\infty$: no penalty on the intercepts), $C' \in \{0.01, 0.03, 0.1\}$, and $\rho_M \in \{1, 0.97, 0.9\}$.
For the experiments with weak labels and the comparison, the A, A~(HuBERT), V, and AV students keep the selected $\gamma$ and $\rho_s$ and reselect $(\rho_M, C', C_b)$ on the same pseudo-targets from $45$ combinations, $\rho_M \in \{1, 0.97, 0.9\}$, $C' \in \{0.003, 0.01, 0.03, 0.1, 0.3\}$, and $C_b \in \{0.5, 2, 8\}$; no label of the target beyond its $K$ draws enters this selection.

\begin{figure}[t]
\centering
\FIG{capture}{\columnwidth}
\caption{Label-free selection does not find the direction error. Fraction of the start's direction error $E_0$ captured by the span of $K=5$ labels per emotion, relative to random selection, for seven label-free rules and for a teaching rule that knows all candidate labels (upper reference), for A, V, and T. Points are ratios of target means; whiskers show the range of the per-target ratios. Ten iterations.}
\label{fig:supp-capture}
\end{figure}

\section{Additional Results for the Comparison}
\label{sec:supp-compare}

\begin{figure*}[t]
\centering
\FIG{compare}{\textwidth}
\caption{Differences from the five-value start at $K=0$ and $K=5$ for the strongest tuned competitor, ours, and ours-ext, per configuration and student (the data of Table~3 of the main text). Bars are means over targets and ten iterations; whiskers show the range across targets. Labels below the bars name the strongest competitor (GD, generalized distillation; KD, distillation with a temperature; HPL, hard pseudo-labels; RB, output rebalancing; ST, self-training; CoT, co-training; CoR, co-regularization).}
\label{fig:supp-compare}
\end{figure*}

\begin{figure}[t]
\centering
\FIG{kcurves}{\columnwidth}
\caption{Differences from the five-value start against $K$ for the four students of the AVT configuration: ours, ours-ext, and the strongest tuned competitor, which is selected again at each $K$.}
\label{fig:supp-kcurves}
\end{figure}

\begin{figure}[t]
\centering
\FIG{entmin}{\columnwidth}
\caption{Differences from the five-value start at $K=0$ for the fused students, per target: fixed entropy minimization and ours-ext. Each point is one target in one configuration.}
\label{fig:supp-entmin}
\end{figure}

\begin{figure}[t]
\centering
\FIG{finetune}{\columnwidth}
\caption{Frozen versus fine-tuned audio encoders. UAR of the frozen five-value start and of our method, and of low-rank adaptation trained on the $K$ labels alone, with the anchor, and with the cross-modal weak labels, with and without centering, at $K=0$ and $K=5$ (four targets, three iterations).}
\label{fig:supp-finetune}
\end{figure}

\section{Fine-Tuning Check}
\label{sec:supp-finetune}

We fine-tuned WavLM with low-rank adapters for the audio student of four targets over three iterations (Fig.~\ref{fig:supp-finetune}).
Fine-tuning on $K=5$ labels per emotion falls \NUM{cmp-ft-konly-below-frozen-e39-k5} points below the frozen five-value start (\NUM{cmp-ft-konly-below-frozen-e39-k5-tgt-lo}--\NUM{cmp-ft-konly-below-frozen-e39-k5-tgt-hi} across targets), and anchoring the adapted head to the start adds \NUM{cmp-ft-anchor-vs-konly-k5} points.
The cross-modal weak labels raise the fine-tuned model by \NUM{cmp-ft-weak-vs-start-k0} points at $K=0$ and \NUM{cmp-ft-weak-vs-anchor-k5} at $K=5$, but the fine-tuned version of our method, with the weak labels of ours-ext fixed to the form selected for the AVT audio student, stays \NUM{cmp-ft-ours-below-frozen-ours-k0} and \NUM{cmp-ft-ours-below-frozen-ours-k5} points below the same configuration on frozen features, below it in all \NUM{cmp-ft-pairs-n} target and iteration pairs at each $K$.
Adding centering raises the fine-tuned start, the model trained on the $K$ labels alone, and the anchored model by \NUM{cmp-ftc-lift-lo}--\NUM{cmp-ftc-lift-hi} points, and the fine-tuned version of our method by \NUM{cmp-ftc-lift-ours-k0}--\NUM{cmp-ftc-lift-ours-k5} points, but it closes only \NUM{cmp-ftc-gap-closed-start-k0}--\NUM{cmp-ftc-gap-closed-konly-k5}\% of the gap to the frozen pipeline.
In this check, the frozen encoder with a linear head is the stronger base with a few labels, and the gain from cross-modal weak labels persists on the fine-tuned encoder.

\section{Full Comparison Tables}
\label{sec:supp-tables}

% 自動生成（scripts/make_supp_tables.py）．手で編集しない．
\begin{table*}[t]
\caption{UAR (\%) of the tuned competitors and of our methods for the single-modality students, averaged over targets and ten iterations. ST: self-training; HPL: hard pseudo-labels; EM: entropy minimization; KD: distillation with a temperature; GD: generalized distillation; CoT: co-training; CoR: co-regularization; PA: prior adjustment by EM; RB: output rebalancing.}
\label{tab:supp-mod}
\centering\footnotesize\setlength{\tabcolsep}{3pt}
\begin{tabular}{@{}llrrrrrrrrrrrr@{}}
\toprule
Config. & Student, $K$ & Start & ST & HPL & EM & KD & GD & CoT & CoR & PA & RB & Ours & Ours-ext \\
\midrule
AVT & A, 0 & \NUM{tab-avt-a-k0-e39} & \NUM{tab-avt-a-k0-self} & \NUM{tab-avt-a-k0-hard} & \NUM{tab-avt-a-k0-entmin} & \NUM{tab-avt-a-k0-kd} & \NUM{tab-avt-a-k0-gd} & \NUM{tab-avt-a-k0-cotrain} & \NUM{tab-avt-a-k0-coreg} & \NUM{tab-avt-a-k0-saerens} & \NUM{tab-avt-a-k0-massbal} & \NUM{tab-avt-a-k0-ours} & \NUM{tab-avt-a-k0-oursx} \\
 & A, 5 & \NUM{tab-avt-a-k5-e39} & \NUM{tab-avt-a-k5-self} & \NUM{tab-avt-a-k5-hard} & \NUM{tab-avt-a-k5-entmin} & \NUM{tab-avt-a-k5-kd} & \NUM{tab-avt-a-k5-gd} & \NUM{tab-avt-a-k5-cotrain} & \NUM{tab-avt-a-k5-coreg} & \NUM{tab-avt-a-k5-saerens} & \NUM{tab-avt-a-k5-massbal} & \NUM{tab-avt-a-k5-ours} & \NUM{tab-avt-a-k5-oursx} \\
AVT & V, 0 & \NUM{tab-avt-v-k0-e39} & \NUM{tab-avt-v-k0-self} & \NUM{tab-avt-v-k0-hard} & \NUM{tab-avt-v-k0-entmin} & \NUM{tab-avt-v-k0-kd} & \NUM{tab-avt-v-k0-gd} & \NUM{tab-avt-v-k0-cotrain} & \NUM{tab-avt-v-k0-coreg} & \NUM{tab-avt-v-k0-saerens} & \NUM{tab-avt-v-k0-massbal} & \NUM{tab-avt-v-k0-ours} & \NUM{tab-avt-v-k0-oursx} \\
 & V, 5 & \NUM{tab-avt-v-k5-e39} & \NUM{tab-avt-v-k5-self} & \NUM{tab-avt-v-k5-hard} & \NUM{tab-avt-v-k5-entmin} & \NUM{tab-avt-v-k5-kd} & \NUM{tab-avt-v-k5-gd} & \NUM{tab-avt-v-k5-cotrain} & \NUM{tab-avt-v-k5-coreg} & \NUM{tab-avt-v-k5-saerens} & \NUM{tab-avt-v-k5-massbal} & \NUM{tab-avt-v-k5-ours} & \NUM{tab-avt-v-k5-oursx} \\
AVT & T, 0 & \NUM{tab-avt-t-k0-e39} & \NUM{tab-avt-t-k0-self} & \NUM{tab-avt-t-k0-hard} & \NUM{tab-avt-t-k0-entmin} & \NUM{tab-avt-t-k0-kd} & \NUM{tab-avt-t-k0-gd} & \NUM{tab-avt-t-k0-cotrain} & \NUM{tab-avt-t-k0-coreg} & \NUM{tab-avt-t-k0-saerens} & \NUM{tab-avt-t-k0-massbal} & \NUM{tab-avt-t-k0-ours} & \NUM{tab-avt-t-k0-oursx} \\
 & T, 5 & \NUM{tab-avt-t-k5-e39} & \NUM{tab-avt-t-k5-self} & \NUM{tab-avt-t-k5-hard} & \NUM{tab-avt-t-k5-entmin} & \NUM{tab-avt-t-k5-kd} & \NUM{tab-avt-t-k5-gd} & \NUM{tab-avt-t-k5-cotrain} & \NUM{tab-avt-t-k5-coreg} & \NUM{tab-avt-t-k5-saerens} & \NUM{tab-avt-t-k5-massbal} & \NUM{tab-avt-t-k5-ours} & \NUM{tab-avt-t-k5-oursx} \\
AVT (2nd) & A, 0 & \NUM{tab-w2-a-k0-e39} & \NUM{tab-w2-a-k0-self} & \NUM{tab-w2-a-k0-hard} & \NUM{tab-w2-a-k0-entmin} & \NUM{tab-w2-a-k0-kd} & \NUM{tab-w2-a-k0-gd} & \NUM{tab-w2-a-k0-cotrain} & \NUM{tab-w2-a-k0-coreg} & \NUM{tab-w2-a-k0-saerens} & \NUM{tab-w2-a-k0-massbal} & \NUM{tab-w2-a-k0-ours} & \NUM{tab-w2-a-k0-oursx} \\
 & A, 5 & \NUM{tab-w2-a-k5-e39} & \NUM{tab-w2-a-k5-self} & \NUM{tab-w2-a-k5-hard} & \NUM{tab-w2-a-k5-entmin} & \NUM{tab-w2-a-k5-kd} & \NUM{tab-w2-a-k5-gd} & \NUM{tab-w2-a-k5-cotrain} & \NUM{tab-w2-a-k5-coreg} & \NUM{tab-w2-a-k5-saerens} & \NUM{tab-w2-a-k5-massbal} & \NUM{tab-w2-a-k5-ours} & \NUM{tab-w2-a-k5-oursx} \\
AVT (2nd) & V, 0 & \NUM{tab-w2-v-k0-e39} & \NUM{tab-w2-v-k0-self} & \NUM{tab-w2-v-k0-hard} & \NUM{tab-w2-v-k0-entmin} & \NUM{tab-w2-v-k0-kd} & \NUM{tab-w2-v-k0-gd} & \NUM{tab-w2-v-k0-cotrain} & \NUM{tab-w2-v-k0-coreg} & \NUM{tab-w2-v-k0-saerens} & \NUM{tab-w2-v-k0-massbal} & \NUM{tab-w2-v-k0-ours} & \NUM{tab-w2-v-k0-oursx} \\
 & V, 5 & \NUM{tab-w2-v-k5-e39} & \NUM{tab-w2-v-k5-self} & \NUM{tab-w2-v-k5-hard} & \NUM{tab-w2-v-k5-entmin} & \NUM{tab-w2-v-k5-kd} & \NUM{tab-w2-v-k5-gd} & \NUM{tab-w2-v-k5-cotrain} & \NUM{tab-w2-v-k5-coreg} & \NUM{tab-w2-v-k5-saerens} & \NUM{tab-w2-v-k5-massbal} & \NUM{tab-w2-v-k5-ours} & \NUM{tab-w2-v-k5-oursx} \\
AVT (2nd) & T, 0 & \NUM{tab-w2-t-k0-e39} & \NUM{tab-w2-t-k0-self} & \NUM{tab-w2-t-k0-hard} & \NUM{tab-w2-t-k0-entmin} & \NUM{tab-w2-t-k0-kd} & \NUM{tab-w2-t-k0-gd} & \NUM{tab-w2-t-k0-cotrain} & \NUM{tab-w2-t-k0-coreg} & \NUM{tab-w2-t-k0-saerens} & \NUM{tab-w2-t-k0-massbal} & \NUM{tab-w2-t-k0-ours} & \NUM{tab-w2-t-k0-oursx} \\
 & T, 5 & \NUM{tab-w2-t-k5-e39} & \NUM{tab-w2-t-k5-self} & \NUM{tab-w2-t-k5-hard} & \NUM{tab-w2-t-k5-entmin} & \NUM{tab-w2-t-k5-kd} & \NUM{tab-w2-t-k5-gd} & \NUM{tab-w2-t-k5-cotrain} & \NUM{tab-w2-t-k5-coreg} & \NUM{tab-w2-t-k5-saerens} & \NUM{tab-w2-t-k5-massbal} & \NUM{tab-w2-t-k5-ours} & \NUM{tab-w2-t-k5-oursx} \\
\bottomrule
\end{tabular}
\end{table*}

\begin{table*}[t]
\caption{UAR (\%) of the tuned competitors and of our methods for the fused students. CoT-F and CoT-P: co-training with the fused or the product view; other abbreviations as in Table~\ref{tab:supp-mod}.}
\label{tab:supp-fused}
\centering\footnotesize\setlength{\tabcolsep}{3pt}
\begin{tabular}{@{}llrrrrrrrrrrr@{}}
\toprule
Config. & $K$ & Start & ST & EM & KD & CoT-F & CoT-P & CoR & PA & RB & Ours & Ours-ext \\
\midrule
AVT & 0 & \NUM{tab-avt-f-k0-e39} & \NUM{tab-avt-f-k0-self} & \NUM{tab-avt-f-k0-entmin} & \NUM{tab-avt-f-k0-kd} & \NUM{tab-avt-f-k0-cotrain-fused} & \NUM{tab-avt-f-k0-cotrain-prod} & \NUM{tab-avt-f-k0-coreg-sum} & \NUM{tab-avt-f-k0-saerens} & \NUM{tab-avt-f-k0-massbal} & \NUM{tab-avt-f-k0-ours} & \NUM{tab-avt-f-k0-oursx} \\
 & 5 & \NUM{tab-avt-f-k5-e39} & \NUM{tab-avt-f-k5-self} & \NUM{tab-avt-f-k5-entmin} & \NUM{tab-avt-f-k5-kd} & \NUM{tab-avt-f-k5-cotrain-fused} & \NUM{tab-avt-f-k5-cotrain-prod} & \NUM{tab-avt-f-k5-coreg-sum} & \NUM{tab-avt-f-k5-saerens} & \NUM{tab-avt-f-k5-massbal} & \NUM{tab-avt-f-k5-ours} & \NUM{tab-avt-f-k5-oursx} \\
AV & 0 & \NUM{tab-av-f-k0-e39} & \NUM{tab-av-f-k0-self} & \NUM{tab-av-f-k0-entmin} & \NUM{tab-av-f-k0-kd} & \NUM{tab-av-f-k0-cotrain-fused} & \NUM{tab-av-f-k0-cotrain-prod} & \NUM{tab-av-f-k0-coreg-sum} & \NUM{tab-av-f-k0-saerens} & \NUM{tab-av-f-k0-massbal} & \NUM{tab-av-f-k0-ours} & \NUM{tab-av-f-k0-oursx} \\
 & 5 & \NUM{tab-av-f-k5-e39} & \NUM{tab-av-f-k5-self} & \NUM{tab-av-f-k5-entmin} & \NUM{tab-av-f-k5-kd} & \NUM{tab-av-f-k5-cotrain-fused} & \NUM{tab-av-f-k5-cotrain-prod} & \NUM{tab-av-f-k5-coreg-sum} & \NUM{tab-av-f-k5-saerens} & \NUM{tab-av-f-k5-massbal} & \NUM{tab-av-f-k5-ours} & \NUM{tab-av-f-k5-oursx} \\
VT & 0 & \NUM{tab-vt-f-k0-e39} & \NUM{tab-vt-f-k0-self} & \NUM{tab-vt-f-k0-entmin} & \NUM{tab-vt-f-k0-kd} & \NUM{tab-vt-f-k0-cotrain-fused} & \NUM{tab-vt-f-k0-cotrain-prod} & \NUM{tab-vt-f-k0-coreg-sum} & \NUM{tab-vt-f-k0-saerens} & \NUM{tab-vt-f-k0-massbal} & \NUM{tab-vt-f-k0-ours} & \NUM{tab-vt-f-k0-oursx} \\
 & 5 & \NUM{tab-vt-f-k5-e39} & \NUM{tab-vt-f-k5-self} & \NUM{tab-vt-f-k5-entmin} & \NUM{tab-vt-f-k5-kd} & \NUM{tab-vt-f-k5-cotrain-fused} & \NUM{tab-vt-f-k5-cotrain-prod} & \NUM{tab-vt-f-k5-coreg-sum} & \NUM{tab-vt-f-k5-saerens} & \NUM{tab-vt-f-k5-massbal} & \NUM{tab-vt-f-k5-ours} & \NUM{tab-vt-f-k5-oursx} \\
AT & 0 & \NUM{tab-at-f-k0-e39} & \NUM{tab-at-f-k0-self} & \NUM{tab-at-f-k0-entmin} & \NUM{tab-at-f-k0-kd} & \NUM{tab-at-f-k0-cotrain-fused} & \NUM{tab-at-f-k0-cotrain-prod} & \NUM{tab-at-f-k0-coreg-sum} & \NUM{tab-at-f-k0-saerens} & \NUM{tab-at-f-k0-massbal} & \NUM{tab-at-f-k0-ours} & \NUM{tab-at-f-k0-oursx} \\
 & 5 & \NUM{tab-at-f-k5-e39} & \NUM{tab-at-f-k5-self} & \NUM{tab-at-f-k5-entmin} & \NUM{tab-at-f-k5-kd} & \NUM{tab-at-f-k5-cotrain-fused} & \NUM{tab-at-f-k5-cotrain-prod} & \NUM{tab-at-f-k5-coreg-sum} & \NUM{tab-at-f-k5-saerens} & \NUM{tab-at-f-k5-massbal} & \NUM{tab-at-f-k5-ours} & \NUM{tab-at-f-k5-oursx} \\
AVT (2nd) & 0 & \NUM{tab-w2-f-k0-e39} & \NUM{tab-w2-f-k0-self} & \NUM{tab-w2-f-k0-entmin} & \NUM{tab-w2-f-k0-kd} & \NUM{tab-w2-f-k0-cotrain-fused} & \NUM{tab-w2-f-k0-cotrain-prod} & \NUM{tab-w2-f-k0-coreg-sum} & \NUM{tab-w2-f-k0-saerens} & \NUM{tab-w2-f-k0-massbal} & \NUM{tab-w2-f-k0-ours} & \NUM{tab-w2-f-k0-oursx} \\
 & 5 & \NUM{tab-w2-f-k5-e39} & \NUM{tab-w2-f-k5-self} & \NUM{tab-w2-f-k5-entmin} & \NUM{tab-w2-f-k5-kd} & \NUM{tab-w2-f-k5-cotrain-fused} & \NUM{tab-w2-f-k5-cotrain-prod} & \NUM{tab-w2-f-k5-coreg-sum} & \NUM{tab-w2-f-k5-saerens} & \NUM{tab-w2-f-k5-massbal} & \NUM{tab-w2-f-k5-ours} & \NUM{tab-w2-f-k5-oursx} \\
\bottomrule
\end{tabular}
\end{table*}

\begin{table}[t]
\caption{Ours-ext minus the strongest tuned competitor over target and iteration pairs: mean (points), effect size $d_z$, Holm-corrected $p$ of the Wilcoxon signed-rank test, and the number of targets with a positive mean.}
\label{tab:supp-stats}
\centering\footnotesize\setlength{\tabcolsep}{2.5pt}
\begin{tabular}{@{}llrrrrrrrr@{}}
\toprule
 & & \multicolumn{4}{c}{$K=0$} & \multicolumn{4}{c}{$K=5$} \\ \cmidrule(lr){3-6}\cmidrule(l){7-10}
Config. & Stud. & Mean & $d_z$ & $p$ & Pos. & Mean & $d_z$ & $p$ & Pos. \\
\midrule
AVT & A & \NUM{st-oursx-best-avt-a-k0-mean} & \NUM{st-oursx-best-avt-a-k0-dz} & $<$0.001 & \NUM{st-oursx-best-avt-a-k0-tgtpos} & \NUM{st-oursx-best-avt-a-k5-mean} & \NUM{st-oursx-best-avt-a-k5-dz} & $<$0.001 & \NUM{st-oursx-best-avt-a-k5-tgtpos} \\
AVT & V & \NUM{st-oursx-best-avt-v-k0-mean} & \NUM{st-oursx-best-avt-v-k0-dz} & \NUM{st-oursx-best-avt-v-k0-pholm} & \NUM{st-oursx-best-avt-v-k0-tgtpos} & \NUM{st-oursx-best-avt-v-k5-mean} & \NUM{st-oursx-best-avt-v-k5-dz} & \NUM{st-oursx-best-avt-v-k5-pholm} & \NUM{st-oursx-best-avt-v-k5-tgtpos} \\
AVT & T & \NUM{st-oursx-best-avt-t-k0-mean} & \NUM{st-oursx-best-avt-t-k0-dz} & \NUM{st-oursx-best-avt-t-k0-pholm} & \NUM{st-oursx-best-avt-t-k0-tgtpos} & \NUM{st-oursx-best-avt-t-k5-mean} & \NUM{st-oursx-best-avt-t-k5-dz} & \NUM{st-oursx-best-avt-t-k5-pholm} & \NUM{st-oursx-best-avt-t-k5-tgtpos} \\
AVT & Fused & \NUM{st-oursx-best-avt-f-k0-mean} & \NUM{st-oursx-best-avt-f-k0-dz} & \NUM{st-oursx-best-avt-f-k0-pholm} & \NUM{st-oursx-best-avt-f-k0-tgtpos} & \NUM{st-oursx-best-avt-f-k5-mean} & \NUM{st-oursx-best-avt-f-k5-dz} & $<$0.001 & \NUM{st-oursx-best-avt-f-k5-tgtpos} \\
AV & Fused & \NUM{st-oursx-best-av-f-k0-mean} & \NUM{st-oursx-best-av-f-k0-dz} & \NUM{st-oursx-best-av-f-k0-pholm} & \NUM{st-oursx-best-av-f-k0-tgtpos} & \NUM{st-oursx-best-av-f-k5-mean} & \NUM{st-oursx-best-av-f-k5-dz} & \NUM{st-oursx-best-av-f-k5-pholm} & \NUM{st-oursx-best-av-f-k5-tgtpos} \\
VT & Fused & \NUM{st-oursx-best-vt-f-k0-mean} & \NUM{st-oursx-best-vt-f-k0-dz} & $<$0.001 & \NUM{st-oursx-best-vt-f-k0-tgtpos} & \NUM{st-oursx-best-vt-f-k5-mean} & \NUM{st-oursx-best-vt-f-k5-dz} & $<$0.001 & \NUM{st-oursx-best-vt-f-k5-tgtpos} \\
AT & Fused & \NUM{st-oursx-best-at-f-k0-mean} & \NUM{st-oursx-best-at-f-k0-dz} & $<$0.001 & \NUM{st-oursx-best-at-f-k0-tgtpos} & \NUM{st-oursx-best-at-f-k5-mean} & \NUM{st-oursx-best-at-f-k5-dz} & $<$0.001 & \NUM{st-oursx-best-at-f-k5-tgtpos} \\
AVT (2nd) & A & \NUM{st-oursx-best-w2-a-k0-mean} & \NUM{st-oursx-best-w2-a-k0-dz} & \NUM{st-oursx-best-w2-a-k0-pholm} & \NUM{st-oursx-best-w2-a-k0-tgtpos} & \NUM{st-oursx-best-w2-a-k5-mean} & \NUM{st-oursx-best-w2-a-k5-dz} & $<$0.001 & \NUM{st-oursx-best-w2-a-k5-tgtpos} \\
AVT (2nd) & V & \NUM{st-oursx-best-w2-v-k0-mean} & \NUM{st-oursx-best-w2-v-k0-dz} & \NUM{st-oursx-best-w2-v-k0-pholm} & \NUM{st-oursx-best-w2-v-k0-tgtpos} & \NUM{st-oursx-best-w2-v-k5-mean} & \NUM{st-oursx-best-w2-v-k5-dz} & \NUM{st-oursx-best-w2-v-k5-pholm} & \NUM{st-oursx-best-w2-v-k5-tgtpos} \\
AVT (2nd) & T & \NUM{st-oursx-best-w2-t-k0-mean} & \NUM{st-oursx-best-w2-t-k0-dz} & \NUM{st-oursx-best-w2-t-k0-pholm} & \NUM{st-oursx-best-w2-t-k0-tgtpos} & \NUM{st-oursx-best-w2-t-k5-mean} & \NUM{st-oursx-best-w2-t-k5-dz} & $<$0.001 & \NUM{st-oursx-best-w2-t-k5-tgtpos} \\
AVT (2nd) & Fused & \NUM{st-oursx-best-w2-f-k0-mean} & \NUM{st-oursx-best-w2-f-k0-dz} & \NUM{st-oursx-best-w2-f-k0-pholm} & \NUM{st-oursx-best-w2-f-k0-tgtpos} & \NUM{st-oursx-best-w2-f-k5-mean} & \NUM{st-oursx-best-w2-f-k5-dz} & $<$0.001 & \NUM{st-oursx-best-w2-f-k5-tgtpos} \\
\bottomrule
\end{tabular}
\end{table}


\bibliographystyle{IEEEtran}
\bibliography{references_ieee}
\end{document}
